\subsection{Puertas cuánticas}
Una puerta de un cúbit $\gate{A}$ actuando sobre un ket $\ket{a_1 a_2}$ puede actuar sobre el primer cúbit, y el resultado es $\gate{A}\ket{a_1}\otimes \ket{a_2}$, o sobre el segundo cúbit, y entonces el resultado es $\ket{a_1}\otimes \gate{A}\ket{a_2}$. Para diferenciar estas operaciones, suele usare la notación tensorial $A\tensor I$ y $I\tensor A$ respectivamente, pero no resulta cómoda cuando el número de cúbits crece y la puerta actúa sobre varios cúbits y en posiciones no consecutivas.

\begin{notation}
  Sea $\gate{A}$ una puerta de un cúbit y $\ket{a}=\ket{a_1\cdots a_n}$ un ket con $n$ cúbits. Dado $\Gamma=\set{k_1,\dots,k_l}\subseteq\N_n$, se usa la siguiente notación
  \[
    \gate{A}_{\Gamma}\ket{a}=
    \gate{A}^{\delta_{1\Gamma}}\ket{a_1}\tensor\cdots\tensor\gate{A}^{\delta_{n\Gamma}}\ket{a_n}
    \,.
  \]
\end{notation}

\begin{remark}
  A lo largo de este trabajo se sigue el convenio estándar de Dirac: las puertas se escriben a la izquierda del ket, y $\gate{A}\ket{a}$ denota la aplicación del operador $\gate{A}$ al estado $\ket{a}$. En consecuencia, en un producto de puertas el factor más a la derecha se aplica primero: $\gate{A}\gate{B}\ket{a}=\gate{A}(\gate{B}\ket{a})$, y un circuito que aplica primero $\gate{B}$ y después $\gate{A}$ se escribe $\gate{A}\gate{B}$, en orden inverso al de su diagrama.
\end{remark}

De esta manera, las dos situaciones comentadas al inicio de esta sección se distinguen escribiendo $\gate{A}_1\ket{a_1 a_2}$ y $\gate{A}_2\ket{a_1 a_2}$, respectivamente.

\begin{proposition}
  Sean $\gate{A}, \gate{B}$ dos puerta de un cúbit.
  Entonces, para cualesquiera dos índices distintos $h,k \in \N$,
  \begin{equation}
    \label{eq:conmutacion-gates}
    \gate{A}_k\gate{B}_{h} = \gate{B}_h\gate{A}_{k}.
  \end{equation}
\end{proposition}
\begin{proof}
  Sea $n\in\N$ tal que $n\geq\max\set{k, h}$ y $\ket{a}=\ket{a_1\cdots a_n}\in\Q_n$ un ket.
  Por definición
  \begin{align*}
    \gate{A}_k\gate{B}_{h}\ket{a}
     & =  \gate{A}_{k}\left(\gate{B}^{\delta_{1h}}\ket{a_1}\tensor\cdots\tensor\gate{B}^{\delta_{nh}}\ket{a_n}\right)                       \\
     & =  \gate{A}^{\delta_{1k}}\gate{B}^{\delta_{1h}}\ket{a_1}\tensor\cdots\tensor\gate{A}^{\delta_{nk}}\gate{B}^{\delta_{nh}}\ket{a_n}\,.
  \end{align*}
  Ahora bien, como $h\neq k$, el producto $\gate{A}^{\delta_{jk}}\gate{B}^{\delta_{jh}}$ para $j=1,\dots,n$ solo puede tomar 3 valores:
  \[
    \gate{A}^{\delta_{jk}}\gate{B}^{\delta_{jh}} = \begin{cases}
      \gate{A} & \text{ si } j = k    \\
      \gate{B} & \text{ si } j = h    \\
      \gate{I} & \text{ en otro caso}
    \end{cases}\,.
  \]
  En todos los casos el producto no depende del orden de los factores, es decir, $\gate{A}^{\delta_{jk}}\gate{B}^{\delta_{jh}}=\gate{B}^{\delta_{jh}}\gate{A}^{\delta_{jk}}$, y llegamos a la igualdad buscada.
\end{proof}

Por la igualdad~(\ref{eq:conmutacion-gates}), la siguiente notación está bien definida.
\begin{notation}
  Sea $\Gamma=\set{k_1,\dots,k_l}\subseteq\mathbb{N}_n$ un conjunto.
  Sea $\gate{A}$ una puerta de un cúbit. Se usa la siguiente notación:
  \begin{align}
    \gate{A}_{\Gamma} & = \gate{A}_{k_1,\dots,k_l}=\prod_{k\in\Gamma}\gate{A}_k\,.
  \end{align}
\end{notation}

Nótese que, de acuerdo con la notación anterior, para el caso particular del conjunto vacío $\Gamma=\emptyset$, la puerta cuántica resultante es la puerta identidad $\gate{A}_\emptyset=\gate{I}$.

Por ejemplo, si $\ket{a}=\ket{a_1 a_2}$ y $\Gamma = \set{1,2}$, entonces
\[
  \gate{A}\ket{a_1}\otimes \gate{A}\ket{a_2} = \gate{A}_{1,2}\ket{a}=\gate{A}_{\Gamma}\ket{a}\,.
\]
